\scanpage{179}
% archive foliation: 178
\uline{Proposition} $S$ un préschéma \uline{connexe},
$S'$ un préschéma sur $S$ tel que $S'\to S$ soit \uline{submersif}
[p.ex. fid plat et quasi-compact, ou propre surjectif].
$S'',S'''$ les carré, cube fibré.

On suppose que $S'$ est \uline{somme} de \uline{préschémas connexes}
(i.e. ses composantes connexes sont ouvertes, i.e. $\pi_0(S')$ discret),
\deleted{et l'ens des \unsure{composantes connexes} de $S''$ est \unsure{q.c.}}

Soit $\xi$ un pt [géométrique] de $S$. Alors pour un groupe discret $G$ variable, le foncteur
\[
G\longmapsto\Pi^1(S'/S,\xi;G)
\]
\deleted{\unsure{considérons}} [où $\Pi^1(S'/S,\xi;G)$ désigne l'ens des classes, à isom près, des données de descente sur $G_{S'}$ \inserted{(fibré principal trivial)} relativement à $S'\to S$, pointées sur $\xi$] est strict$^{\mathrm t}$ proreprésentable :
\[
\Pi^1(S'/S,\xi;G)=\varinjlim\operatorname{Hom}(\pi_i,G)
\]

Si $\pi_0(S'')$ est \inserted{quasi-}\uline{compact}, i.e. tte fonction loc constante sur $\pi_0(S'')$ \unsure{n'a qu'un nb fini de valeurs}, --- p.ex. $S''$ quasi-compact]
\deleted{\unsure{ou $S''$ a}} \deleted{[un nb fini de comp. de connexité]}
\uline{alors} les $\pi_i$ sont de type fini. Si $S''$ a un nb fini \inserted{$n''$} de composantes connexes, alors le foncteur précédent est représentable : \deleted{de façon} \deleted{par}
\[
\Pi^1(S'/S,\xi;G)\simeq\operatorname{Hom}(\pi_1(S'/S,\xi),G)
\]
avec $\pi_1(S'/S,\xi)$ \deleted{\unsure{groupe}} \inserted{groupe discret} de type fini.
De façon précise, * Si \unsure{le nb} $n'$ des comp. conn. de $S'$ est aussi fini, et \deleted{$\unsure{n''-n'}$} $\pi_1(S'/S;\xi)$ est engendré par $\unsure{n''-n'}$ générateurs.
